% !TeX program = pdflatex
% Self-contained source. Compile twice with pdflatex.
\documentclass[11pt,a4paper]{article}
\usepackage[T1]{fontenc}
\usepackage[utf8]{inputenc}
\usepackage{lmodern,microtype,amsmath,amssymb,amsthm,booktabs,tabularx}
\usepackage[margin=25mm,top=24mm,bottom=25mm,headheight=14pt]{geometry}
\usepackage{xcolor,fancyhdr,listings,enumitem}
\usepackage[colorlinks=true,linkcolor=black!55!teal,citecolor=black!55!teal,urlcolor=black!55!teal]{hyperref}
\definecolor{forest}{HTML}{165154}
\definecolor{softgray}{HTML}{F4F6F5}
\newtheorem{theorem}{Theorem}[section]
\newtheorem{lemma}[theorem]{Lemma}
\newtheorem{proposition}[theorem]{Proposition}
\newtheorem{corollary}[theorem]{Corollary}
\theoremstyle{definition}\newtheorem{problem}[theorem]{Problem}
\renewcommand{\sectionautorefname}{Section}
\setlength{\emergencystretch}{2em}
\setlength{\parindent}{1em}
\setlength{\parskip}{2pt}
\raggedbottom
\setlist{topsep=3pt,itemsep=2pt}
\pagestyle{fancy}\fancyhf{}
\fancyhead[L]{\sffamily\small Degree differences and log-concavity}
\fancyhead[R]{\sffamily\small Research manuscript}
\fancyfoot[L]{\sffamily\footnotesize 15 November 2026}
\fancyfoot[R]{\sffamily\small\thepage}
\renewcommand{\headrulewidth}{0.3pt}
\lstset{basicstyle=\ttfamily\footnotesize,columns=fullflexible,keepspaces=true,showstringspaces=false,breaklines=true,backgroundcolor=\color{softgray},frame=single,framerule=0pt,framesep=5pt,aboveskip=9pt,belowskip=9pt}
\hypersetup{pdftitle={Antiregular graphs maximize every increasing concave degree-difference sum},pdfsubject={Extremal graph theory, chromatic polynomials and binomial inequalities}}
\begin{document}
\thispagestyle{plain}
\noindent{\sffamily\small\bfseries\color{forest} RESEARCH MANUSCRIPT\hfill 15 NOVEMBER 2026}\par
\vspace{9pt}
\noindent{\sffamily\bfseries\fontsize{24}{28}\raggedright\selectfont Antiregular graphs maximize every \mbox{increasing} concave degree-difference sum\par}
\vspace{7pt}
\noindent{\large Extremal chemical graphs, iterated log-concavity and two binomial sequences}\par
\vspace{11pt}
\noindent\textbf{Abstract.}
Among simple graphs of a given order, the same antiregular graph maximizes
the sum of any nondecreasing discretely concave function of the absolute
difference of two vertex degrees. A supporting line chosen by the difference
of the two \emph{ranks} gives the result and an optimal stability bound.
This proves Problems 5 and 6 of Knor, \v{S}krekovski, Filipovski and
Dimitrov. A separate integer inequality settles their Conjecture 10 for
chemical graphs, including the three residue classes outside the scope
of the authors' 2026 follow-up.
We also disprove Amdeberhan and Moll's Conjecture 21 and classify the
cycle lengths with infinitely log-concave chromatic coefficients.
For two binomial sums we prove Brian Yi Sun's Conjecture 4.7,
give an opposite-sign theorem for his Conjecture 4.8, and prove
Zhi-Wei Sun's scalar log-convexity conjecture for all its exponents.
The algebraic certificates use exact arithmetic. The main graph
inequalities are proved without computation.
\vspace{5pt}
\section*{Original questions and outcomes}
\small
\renewcommand{\arraystretch}{1.16}
\begin{tabularx}{\linewidth}{@{}p{.30\linewidth}X@{}}
\toprule
\textbf{Source} & \textbf{Outcome established here}\\
\midrule
Knor et al., Problems 5, 6 & The two antiregular graphs maximize at $p=1/n$ and at every $0<p<1$. They are the only maximizers up to isomorphism.\\
Knor et al., Conjecture 10 & Exact degree multiplicities for every connected chemical graph extremizer of every order $n\ge7$ at $p=1/n$.\\
Amdeberhan--Moll, Conj. 21 & False. Every cycle of order at least $17$ fails by the third iterate. With zero extension the full cutoff for infinite log-concavity is $n\le11$.\\
Brian Yi Sun, Conjs. 4.7, 4.8 & The successive ratios satisfy the first conjecture strictly. The second fails at its first interior index and satisfies the opposite inequality throughout its tail.\\
Zhi-Wei Sun, scalar conjecture & Strict log-convexity for all required $m\ge5$ at $q=1$, inside a larger region of real parameters.\\
\bottomrule
\end{tabularx}
\normalsize
\vspace{6pt}
\noindent The source and priority comparison is in Section~\ref{sec:status}.
The two Sun entries concern different authors and different sequences.
\clearpage
\section{A common maximizer for nondecreasing concave weights}
\label{sec:general}
Label the vertices of a finite simple undirected graph $G$ so that
$d_1\le\cdots\le d_n$. For $f\colon\{0,\ldots,n-1\}\to\mathbb R$ set
\[
 I_f(G)=\sum_{i<j} f(d_j-d_i),\qquad I_p(G)=\sum_{i<j}(d_j-d_i)^p
 \quad(p>0).
\]
Write $A_n$ for the graph on $1,\ldots,n$ having edge $ij$ ($i<j$)
exactly when $i+j\ge n+1$. Its increasing degree list is
\begin{equation}\label{eq:rankgap}
 d_i^*=i-\mathbf1_{\{2i\ge n+1\}},\qquad
 d_j^*-d_i^*\in\{j-i-1,j-i\}.
\end{equation}
Indeed, the potential neighbors of $i$ are the $i$ ranks at least $n+1-i$,
and its own rank must be excluded if it lies among them. There are $n-1$
distinct degrees, with exactly one repetition. Such a graph is called
antiregular.

\paragraph{The auxiliary question.}
Can one choose a linear upper bound for the contribution of each ranked
pair before any edges of the competing graph are known, and have every
bound be attained on the same graph? This asks for more than
an optimizer for a particular power. For an arbitrary nondecreasing
discretely concave $f$, the following choice answers the question.

\begin{lemma}[Support by rank difference]\label{lem:rank}
Set $s_r=f(r)-f(r-1)$ for $1\le r\le n-1$, with
$s_1\ge\cdots\ge s_{n-1}\ge0$. Put $b_r=f(r)-r s_r$ and
$g_i=f(i-1)-f(n-i)$. Then
\[
 I_f(G)\le \sum_{i<j} b_{j-i}
                +\sum_{\{i,j\}\in E(G)}(g_i+g_j).
\]
The graph $A_n$ attains equality in this bound and also maximizes its
right-hand side over all simple graphs with ranks $1,\ldots,n$.
\end{lemma}
\begin{proof}
For an integer $t$ in the domain,
\[
 f(t)\le b_r+s_r t .
\]
At $r-1$ and $r$ this is equality. To the left sum slope increments
at least $s_r$, and to the right sum increments at most $s_r$.
Use the inequality at $r=j-i$, $t=d_j-d_i$. By
\eqref{eq:rankgap} all inequalities are equalities for $A_n$.
The coefficient of $d_i$ in the linear part is
\[
  \sum_{r=1}^{i-1}s_r-\sum_{r=1}^{n-i}s_r
   = f(i-1)-f(n-i)=g_i .
\]
Counting edge endpoints rewrites the linear part as the displayed edge
sum. The sequence $g_i$ is nondecreasing and $g_{n+1-i}=-g_i$.
Consequently $g_i+g_j$ is nonnegative for $i+j\ge n+1$ and nonpositive
for $i+j\le n+1$. The graph $A_n$ includes every positive-weight edge
and no negative-weight edge. This proves the last assertion, including
rank ties and zero-weight edges.
\end{proof}

\begin{theorem}[Extremality and sharp stability]\label{thm:rank}
Let $n\ge4$. Let $f$ be nondecreasing and discretely concave, and put
$\kappa=2f(1)-f(0)-f(2)$. If $x_d$ denotes the multiplicity of degree
$d$ in $G$, then
\begin{equation}\label{eq:stability}
 I_f(A_n)-I_f(G)\ \ge \kappa\left(\sum_d\binom{x_d}{2}-1\right).
\end{equation}
When $\kappa>0$, equality of the two scores occurs precisely for $A_n$
and its complement, up to isomorphism. The coefficient in
\eqref{eq:stability} is optimal uniformly in the graph order.
\end{theorem}
\begin{proof}
The function $u(t)=\min(t,1)$ counts pairs of unequal degrees.
Subtracting $\kappa u$ from $f$ changes only the first slope, replacing
it by $s_2$. The remaining function is still nondecreasing and
discretely concave, so Lemma~\ref{lem:rank} applies. The difference
between the two $I_u$ scores is exactly the bracket in
\eqref{eq:stability}. Some degree must repeat, since a vertex of degree
zero and one of degree $n-1$ cannot coexist. If scores tie and
$\kappa>0$, precisely one pair has equal degrees.

For completeness the last condition characterizes the two claimed
graphs. If a graph with one repeated pair has an isolated vertex, pass
to its complement, which has none and has the same score. All degrees
$1,\ldots,n-1$ now occur. For $n\ge4$ its leaf forces its universal
vertex to be unique. Two leaves would exclude degree $n-2$, since a
vertex of that degree would miss both leaves. Thus the leaf is unique.
Deleting the leaf and the universal vertex reduces every remaining
degree by one. The remaining graph has all degrees $1,\ldots,n-3$.
Induction with base graphs $K_2,P_3$ and reversal of this construction
gives $A_n$, which has the stated degrees by \eqref{eq:rankgap}.
The two graphs attain the same score by invariance of absolute
degree differences.

For sharpness compare $A_4$ (a triangle with one pendant vertex) with $K_{1,3}$. Their counts at
differences $(0,1,2)$ are $(1,4,1)$ and $(3,0,3)$. The left side
of \eqref{eq:stability} is $2\kappa$ and the bracket equals $2$.
\end{proof}

\begin{corollary}[Problems 5 and 6 of Knor et al.]\label{cor:powers}
For each $n\ge4$ and every $0<p<1$ the only maximizers of $I_p$
are the two antiregular graphs. This includes $p=1/n$ for every $n$.
The stability coefficient is $2-2^p$.
\end{corollary}
\begin{proof}
The successive differences of $t^p$ are positive and decreasing and
$2-2^p>0$. Apply Theorem~\ref{thm:rank}.
\end{proof}
For $p=1$ the antiregular graph is a maximizer, as already proved
with the full equality classification by Abdo, Brandt and Dimitrov
\cite{abd}. The star calculation above gives a tie when $p=1$ and
a strict reversal for every $p>1$. Thus among positive powers, $(0,1]$
is the exact range for which $A_n$ maximizes at \emph{every} order.
A useful equivalent formulation of the common maximization result is
that it holds for each cap $f_k(t)=\min(t,k)$. These functions generate
the nondecreasing concave functions, up to constants, by summing
successive slope differences. The argument uses no asymptotic passage
as $p$ tends to zero.
\section{The chemical graph conjecture at exponent \texorpdfstring{$1/n$}{1/n}}
\label{sec:chem}

Knor, \v{S}krekovski, Filipovski and Dimitrov
\cite[Theorem~9 and Conjecture~10]{ksfd} prescribe the multiplicities of
degrees $1,2,3,4$ maximizing
\[
 I_p(G)=\sum_{\{u,v\}\subseteq V(G)}|d(u)-d(v)|^p
\]
among connected chemical graphs. A chemical graph is connected and simple and has maximum
degree at most four. Their conjecture takes $p=1/n$ for $n\ge7$.
The following argument gives a sufficient interval of exponents.

\begin{theorem}\label{thm:chem}
Suppose $n\ge7$, $p>0$ and $3^p-1<4/(3n)$. Write $n=4k+r$, with
$0\le r\le3$. The maximizing degree multiplicity vectors are exactly
$k(1,1,1,1)+s$ for the following offsets:
\[
\begin{array}{c|l}
r&s\\ \hline
0&(0,0,0,0)\\
1&(0,0,0,1)\\
2&(0,1,0,1),\ (1,0,1,0)\\
3&(1,0,1,1).
\end{array}
\]
In particular the conclusion holds at $p=1/n$.
\end{theorem}

\begin{proof}
Write $x=(x_1,x_2,x_3,x_4)$ for a nonnegative integer multiplicity vector
with total $n$. First allow all such vectors, without imposing graphicality.
Its objective is
\[
 I_p(x)=\tfrac12(n^2-\textstyle\sum_i x_i^2)+H(x),\qquad
 H(x)=A(x_1x_3+x_2x_4)+B x_1x_4,
 \quad A=2^p-1,\quad B=3^p-1.
\]
Here $0<A<B$. Put $k=\lfloor n/4\rfloor$ and write $x=k\mathbf{1}+z$, so
$\sum_i z_i=r$. Fix one displayed vector $y=k\mathbf{1}+s$ of the same
residue class. Let
\[
 Q=\sum_i z_i^2,\quad D=(Q-r)/2,\quad E(z)=z_1+z_4,qquad
 H(z)=A(z_1z_3+z_2z_4)+Bz_1z_4.
\]
The identical quadratic notation $H$ applies also to signed offsets.
Since $j(j-1)$ is a nonnegative even integer for every integer $j$,
$D\ge0$ is integral, and $D=0$ precisely when each $z_i$ is zero or
one. An exact expansion gives
\begin{equation}\label{eq:chem-margin}
 I_p(x)-I_p(y)=-D+kB(E(z)-E(s))+H(z)-H(s).
\end{equation}
We consider only the necessary parity condition $z_1+z_3$ even.

First, $E(z)-E(s)\le D$. To see this, subtract one from coordinates
1 and 4 of $z$. The resulting integer vector $w$ has sum $r-2$,
$w_1+w_3$ odd, and squared norm $Q-2E(z)+2$. The minimum possible norm
squared is at least $2,1,2,1$ for $r=0,1,2,3$, respectively. In the
first, second, and fourth cases use $\sum_i w_i^2\ge|\sum_i w_i|$.
In the third case a zero-sum vector with the required odd parity is
nonzero and therefore has squared norm at least two.
Translating $s$ in the same way attains the stated lower bound in every
case. Comparing its squared norm with that of $w$ gives the desired
inequality.

For every signed $z$,
$H(z)\le BQ$ by $2|uv|\le u^2+v^2$ and $A\le B$; each coordinate
belongs to at most two of the three pairs. Also $H(s)\ge0$. For
$D\ge1$ and $n\ge10$, this gives
\[
 H(z)-H(s)\le B(r+2D)\le (r+2)BD\le (n/2)BD.
\]
When $n=8$ or $9$, the last inequality also holds since $r=0$ or $1$.
When $n=7$, $r=3$ and $H(s)=A+B\ge B$, so that
$H(z)-H(s)\le(2+2D)B\le4BD$.
Consequently for $n\ge8$ the positive terms in
\eqref{eq:chem-margin} are bounded above by
$(k+n/2)BD\le(3n/4)BD<D$. For $n=7$ the bound is $(1+4)BD<D$,
since $B<4/21$. Every unbalanced admissible vector is thus strictly
worse than $y$.

When $D=0$, the coordinates of $z$ are zeros and ones. In residue $0$
there is only one choice. In residue $1$ parity allows the additional
unit at degree $2$ or $4$, and the second choice wins by $kB>0$.
In residue $3$ parity allows omission of degree $2$ or $4$. Both
offsets have the $A$-pair $(1,3)$; omission of $2$ gains $(k+1)B$.
In residue $2$ parity leaves precisely $\{1,3\}$ or $\{2,4\}$ for
the two extra units. They give equal values by direct substitution.

It remains to realize the degree lists. The original
proof of \cite[Theorem~9]{ksfd} verifies their graphicality for $n\ge7$.
All entries are positive, and each displayed list has degree sum at least
$2(n-1)$ (substitution gives $10k,10k+4,10k+4$ or $10k+6$,
and $10k+8$ in the four residues).
Among all simple realizations take one with the fewest components.
If it has two or more, some component contains a cycle, since a forest
with at least two components has at most $n-2$ edges. Remove a cycle edge
$uv$ there and any edge $ab$ in another component, and insert $ua,vb$.
This preserves every degree and joins the two components, proving
connected realizability.

Finally $(3^t-1)/t$ is increasing for $t>0$ by convexity of $3^t$.
Hence $n(3^{1/n}-1)\le7(3^{1/7}-1)<4/3$ for $n\ge7$.
The final strict inequality follows from
$(1+4/21)^7\ge1+7(4/21)+21(4/21)^2>3$.
\end{proof}
\section{Which cycles have infinitely log-concave chromatic coefficients?}
\label{sec:cycles}
For a finite coefficient sequence $a$, extend it by zero and put
$(\mathcal La)_k=a_k^2-a_{k-1}a_{k+1}$. Infinite log-concavity requires
every entry of every positive iterate to be nonnegative. Huh proved
the first-iterate inequality for the absolute coefficients of every
chromatic polynomial \cite{huh}. Amdeberhan and Moll asked whether it
persists at every iterate \cite[Conjecture 21]{am}. The same question
is Problem~5 in the report of the 2022 BIRS workshop \cite{birs}.

\begin{theorem}\label{thm:cycle}
For $n\ge3$, the sequence of absolute coefficients of the chromatic
polynomial of $C_n$ is infinitely log-concave if and only if $n\le11$.
For every $n\ge17$ its third iterate is negative at degree $2$.
In particular Amdeberhan and Moll's conjecture is false.
\end{theorem}
\begin{proof}
The chromatic polynomial of a cycle is
$(q-1)^n+(-1)^n(q-1)$ \cite{am}. Hence $a_0=0$, $a_1=n-1$, and
$a_k=\binom nk$ for $2\le k\le n$. Only five nonconstant
coefficients can affect the third iterate at index two. Applying
$\mathcal L$ twice to these gives
\begin{align*}
 (\mathcal L^2a)_1&=(n-1)^4,\\
 (\mathcal L^2a)_2&=n^2(n-1)^4(n+2)/16,\\
 (\mathcal L^2a)_3&=n^3(n-2)^2(n-1)^4(n+1)(n^2+n+18)/51840 .
\end{align*}
The next subtraction factors as
\begin{align}
 (\mathcal L^3a)_2&=-\frac{n^3(n-1)^8(n+4)}{103680}\,P(n),\label{eq:cycle}\\
 P(n)&=2n^4-12n^3-327n^2-412n+36.\nonumber
\end{align}
For $n=17+h$ this quartic is
$2h^4+124h^3+2529h^2+17370h+6615$, positive for all $h\ge0$.
A witness small enough to check by hand is
\[
 (\mathcal L^3 a(C_{17}))_2
 =22491136^2-65536\cdot8150077440
 =-28272276537344 .
\]
If iteration is defined only at strictly interior indices and both
endpoints are discarded after each step, take the graph
$C_{17}\sqcup 2K_1$. Its chromatic polynomial is shifted by $q^2$.
The same witness is at index four and depends only on original
indices one through seven, all interior. This gives a counterexample
under that interpretation as well.

To classify the smaller cycles, let $\rho=(3+\sqrt5)/2$.
If a nonnegative sequence $a$ satisfies $a_i^2\ge\rho a_{i-1}a_{i+1}$
throughout, then $b=\mathcal La$ satisfies it as well
\cite[Lemma 2.1]{mcs}. Indeed
$(1-1/\rho)a_i^2\le b_i\le a_i^2$ and $(\rho-1)^2=\rho$.
For each short cycle the table gives either an iterate in this
invariant region (with all earlier iterates nonnegative), or an
iterate with a negative entry at index two. The displayed integer
checks verify every entry needed for the table.
\[
\begin{array}{c|ccccc|cc}
 n&3,4&5&6,7,8,9,10&11&&12&13,14,15,16\\ \hline
 \text{iterate}&0&1&2&3&&5&4\\
 \text{test}&\multicolumn{4}{c}{\text{invariant region}}&&
 \multicolumn{2}{c}{\text{negative at index two}}
\end{array}
\]
\begin{lstlisting}[language=Python]
from math import comb
def step(a):
    b = [0] + a + [0]
    return [y*y-x*z for x,y,z in zip(b,b[1:],b[2:])]
def in_region(a):
    for x,y,z in zip(a,a[1:],a[2:]):
        v = 2*y*y - 3*x*z
        if v < 0 or v*v < 5*(x*z)**2:
            return False
    return min(a) >= 0
for n, depth in zip(range(3,17), (0,0,1,2,2,2,2,2,3,5,4,4,4,4)):
    a = [0] + [comb(n,k)-int(k==1) for k in range(1,n+1)]
    for t in range(depth):
        assert min(a) >= 0
        a = step(a)
    if n <= 11:
        assert in_region(a)
    else:
        assert a[2] < 0
\end{lstlisting}
The invariant-region implication covers arbitrarily many future
iterations for $n\le11$. Equation~\eqref{eq:cycle} covers the
unbounded range of lengths.
\end{proof}

\subsection{A local obstruction that survives attaching a clique}
The exact quartic also has a structural explanation. Suppose $a^{(n)}$
are sequences extended by zero, with zero constant term. If for some
$c>0$ and positive $s_n$ one has
$a^{(n)}_j/s_n^j\to c/j!$ for $1\le j\le5$, then
\[
 s_n^{-16}(\mathcal L^3a^{(n)})_2\longrightarrow-c^8/51840.
\]
Indeed, the operation is local and polynomial. Multiplying $a_j$ by
$s^j$ multiplies its $t$th iterate at $j$ by $s^{2^t j}$.
For the limiting input $u_0=0$, $u_j=1/j!$ for $1\le j\le5$,
the entries at $1,2,3$ of $\mathcal L^2u$ are $1,0,1/51840$.
The vanishing central entry leaves a strictly negative third iterate.
This provides an obstruction for a neighborhood of the rescaled
coefficient profile.

For instance, fix $m\ge3$ and identify an edge of $C_n$ with
an edge of $K_m$. Once the colors of the two shared vertices
are given, the remaining clique vertices can be colored in
$(q-2)\cdots(q-m+1)$ ways. Alternating signs turn multiplication
by this factor into convolution of the absolute coefficients with
those of $\prod_{j=2}^{m-1}(z+j)$. These coefficients are fixed
and its constant term is $(m-1)!$. At indices $1\le j\le5$,
division by $n^j$ leaves the limit $(m-1)!/j!$.
Thus the third iterate fails for all sufficiently large $n$ for
each $m$. The two graphs being glued remain connected upon deleting
any one vertex, and they retain a shared vertex after any single
deletion. The resulting graphs are therefore 2-connected. They have
chromatic number $m$ because $m\ge3$ and a proper cycle coloring
extends across the shared edge to a coloring of the clique.
\section{A simultaneous certificate for Sun's two conjectures}
\label{sec:sun}
For $n\ge0$ let
\[
 R_n=\sum_{k=0}^n\binom nk\binom{n+k}{k}\frac1{2k-1},
 \qquad r_n=R_{n+1}/R_n\quad(n\ge1).
\]
Here $R_n>0$ for $n\ge1$, since the terms at $k=0,1$ sum to
$n(n+1)-1>0$, and the remaining summands are positive.
In the 2015 preprint and its 2018 published version, Brian Yi Sun
asks in Conjecture 4.7 whether $(r_n)_{n\ge4}$ is log-concave.
Conjecture 4.8 asks whether $(R_n^2-R_{n-1}R_{n+1})_{n\ge6}$
is infinitely log-concave \cite{bysun}. These questions follow the main
theorem of that paper, which establishes monotonicity of $r_n$ and a
separate inequality for roots of $R_n$.
\begin{theorem}\label{thm:sun-ratio}
For all $n\ge5$ one has $r_n^2>r_{n-1}r_{n+1}$.
At $n=4$ the inequality reverses. In particular, Conjecture 4.7
holds strictly at every interior index of its stated sequence.
\end{theorem}
\begin{theorem}\label{thm:sun-minor}
Put $V_n=R_{n-1}R_{n+1}-R_n^2$. For every $n\ge6$,
\[
                V_{n-1}V_{n+1}>V_n^2 .
\]
The inequality reverses at $n=5$.
Conjecture 4.8 as stated is false, with failure of its first required
log-concavity inequality already at the interior index $n=7$.
\end{theorem}
Indeed, letting $D_n=-V_n$ gives
$(D_6,D_7,D_8)=(-127448,-1712480,-25704608)$ and
$D_7^2-D_6D_8=-343413129984$. This test is independent of the convention
for the first endpoint of a one-sided sequence.

\subsection{The auxiliary enclosure problem}
We seek one pair of bounds narrow enough to decide independently both
the ratio curvature and the curvature of its first determinant.
For the recurrence normalization put
$t=3-2\sqrt2$ and $\lambda=t^{-1}=3+2\sqrt2$.
Generating the binomial sum, as a formal power series, gives
\[
 G(z)=\sum_{n\ge0}R_nz^n
     =-\frac{\sqrt{1-6z+z^2}}{(1-z)^2}.
\]
In fact summation over $n\ge k$ at fixed $k$ yields
$\binom{2k}k z^k/(1-z)^{2k+1}$ before the factor $(2k-1)^{-1}$.
The remaining sum follows from
$\sum_{k\ge0}\binom{2k}k w^k/(2k-1)=-\sqrt{1-4w}$.
Differentiation gives
$(1-z)(1-6z+z^2)G'=(z^2-8z-1)G$, whence
\[
 (n+3)R_{n+3}=(7n+13)R_{n+2}-(7n+15)R_{n+1}+(n+1)R_n.
\]
The starting values are $(-1,1,7)$. Thus $y_n=t r_n$ obeys
$y_{n+2}=F_n(y_n,y_{n+1})$, where
\[
 F_x(a,v)=\frac{(7x+13)t-(7x+15)t^2/v+(x+1)t^3/(av)}{x+3}.
\]
For $x\ge24$, $a>1/2$, $v>0$, the map decreases in $a$ and
increases in $v$. For the second assertion use
$(7x+15)a>(x+1)t$.

Define the constants and two polynomials
\[
 b=\frac{84-39\sqrt2}{32},\qquad
 c=\frac{-537+273\sqrt2}{64},\qquad
 d=\frac{108996-83715\sqrt2}{4096},
\]
\[
 H_\pm(x)=x^5-\tfrac32x^4+b x^3+c x^2+d x\pm20.
\]
One can derive the center by substituting
$S(n)=1+\sum_{j=1}^4a_j n^{-j}$ into
\begin{align*}
 &(1+3/n) S(n)S(n+1)S(n+2)-(7+13/n)t S(n)S(n+1)\\
 &\hspace{20mm}+(7+15/n)t^2 S(n)-(1+1/n)t^3 .
\end{align*}
Setting the coefficients of $n^{-1},\ldots,n^{-4}$ to zero gives
$(-3/2,b,c,d)$ successively. The coefficient of the new unknown at
each step is $3-14t+7t^2>0$. The proof below establishes the error
bound from the recurrence itself. The use of ratio enclosures has
precedents in both Sun's interlacing argument and the work of
Hou and Zhang on $r$-log-convexity \cite{bysun,houzhang}.

\begin{lemma}\label{lem:sun-tube}
For every $n\ge24$, one has
$H_-(n)/n^5<y_n<H_+(n)/n^5$.
\end{lemma}
\begin{proof}
Put $x=24+h$ and write $d_i=(x+i)^5$, $l_i=H_-(x+i)$,
$u_i=H_+(x+i)$. For $a,b',c'$ in the polynomial ring
$\mathbb{Q}(\sqrt2)[h]$ define
\begin{align*}
 P(a,b',c')={}&(7x+13)t ab'd_2-(7x+15)t^2 a d_1d_2\\
               &+(x+1)t^3d_0d_1d_2-(x+3)ab'c'.
\end{align*}
The remaining polynomial certificates are
\begin{align*}
 A&=l_1d_0-u_0d_1,& B&=l_3d_2-u_2d_3,& E&=u_2d_1-l_1d_2,\\
 Q&=l_1^2d_0d_2-u_0u_2d_1^2,&
 T&=l_2ABd_1-u_0 E^2d_3 .
\end{align*}
The exact expansion checks in Appendix~\ref{app:ratio} give strictly
positive coefficients in every degree from zero through the degree shown:
\[
\begin{array}{c|rrrrrrr}
 \text{polynomial}& l_0-d_0/2&P(u_0,l_1,l_2)&-P(l_0,u_1,u_2)& Q&A&B&T\\
 \hline
 \text{degree}&5&11&11&17&8&8&24
\end{array}
\]
Each polynomial is positive for every real $h\ge0$. In particular,
$l_i/d_i>1/2$. By clearing positive denominators, the two signs of $P$
give
\[
 F_x(u_0/d_0,l_1/d_1)>l_2/d_2,\qquad
 F_x(l_0/d_0,u_1/d_1)<u_2/d_2.
\]
The initial tests at indices $24,25$ use respectively
$R_{24}=8492182185653327$, $R_{25}=46471113779766769$ and
$R_{26}=254949930551822839$. Monotonicity of $F_x$ on this
rectangle then proves the assertion by induction.
\end{proof}

\begin{proof}[Proof of Theorem~\ref{thm:sun-ratio}]
The sign of $Q$ gives
$(l_1/d_1)^2>(u_0/d_0)(u_2/d_2)$ for $x\ge24$. The lemma proves the
assertion for centers $n=x+1\ge25$.
For $5\le n\le24$ the recurrence gives
$R_{n+1}^3R_{n-1}-R_n^3R_{n+2}>0$.
At $n=4$ this integer equals $-4623352$.
Appendix~\ref{app:ratio} gives the finite arithmetic in full.
\end{proof}
\begin{proof}[Proof of Theorem~\ref{thm:sun-minor}]
Write $V_n=\lambda R_{n-1}R_n(y_n-y_{n-1})$.
On cancelling positive factors, the sign of $V_{n-1}V_{n+1}-V_n^2$
is the sign of
\[
 y_n(y_{n-1}-y_{n-2})(y_{n+1}-y_n)
               -y_{n-2}(y_n-y_{n-1})^2 .
\]
Take $x=n-2\ge24$. The positive numerators $A$ and $B$ bound the two factors in the first product from below by
$A/(d_0d_1)$ and $B/(d_2d_3)$. Applying the same $A$ bound at
$x+1$ shows that $0<y_n-y_{n-1}<E/(d_1d_2)$.
The entire expression therefore exceeds
$T/(d_0d_1^2d_2^2d_3)>0$. Direct recurrence evaluation gives the
same sign for $6\le n\le25$, and gives $-16234176$ at $n=5$.
The explicit test at $n=7$ above concludes the disproof.
\end{proof}
\section{A scalar conjecture of Zhi-Wei Sun}
\label{sec:sunbin}

Set
\[
 S_n^{(m)}(q)=\sum_{k=0}^n\binom{n+k}{k}^{m}q^k,\qquad n\ge0.
\]
On 25 July 2026 Sun conjectured that $(S_n^{(m)}(1))_{n\ge0}$
is log-convex for all integers $m\ge5$, in a comment following Zhou's
counterexample to the earlier coefficientwise assertion \cite{zwsun}.
We prove strict inequalities on a continuous set of parameters.

\begin{theorem}\label{thm:sunbin}
If $m\ge1$ is real and $q\ge 2^{2-m}$, then for every $n\ge1$,
\[
 S_{n-1}^{(m)}(q) S_{n+1}^{(m)}(q)>\bigl(S_n^{(m)}(q)\bigr)^2 .
\]
In particular Sun's scalar conjecture holds for every specified exponent.
\end{theorem}

Our auxiliary problem asks how much the sum of the terms measured
relative to the last term can increase in one step. The estimate below
combines a product telescope with a differentiated geometric series.
These tools also occur in the counting comparison and the weighted
length sums of VMRC manuscript [043] \cite{vmrc}. Here the factor
difference must retain its order $n^{-2}$ to compete with the curvature
of the central binomial coefficient.

\begin{lemma}\label{lem:sunbin-drift}
Let $B_n=\binom{2n}{n}$, $a=2^{-m}/q$ and
$H_n=S_n^{(m)}(q)/(q^n B_n^m)$, with $m\ge1$ and $0<a<1$.
Then $H_n$ is nondecreasing and for $n\ge2$,
\[
 0\le\frac{H_n}{H_{n-1}}-1
 \le \frac{2ma^2}{(1-a)^3(1+a)(n+1)(n-1)}.
\]
\end{lemma}
\begin{proof}
Reversing the summation index with $j=n-k$ gives
\[
 H_n=\sum_{j=0}^n a^j
     \left(\prod_{i=0}^{j-1}u_i(n)\right)^m,\qquad
 u_i(n)=\frac{2(n-i)}{2n-i}=1-\frac{i}{2n-i}.
\]
For $n\ge2$ compare the sums at $n$ and $n-1$, including a zero term
of index $j=n$ at $n-1$ by using $u_{n-1}(n-1)=0$. For $0\le i<j\le n$
all factors lie in $[0,1]$, and
\[
 0\le u_i(n)-u_i(n-1)
 =\frac{2i}{(2n-i)(2n-2-i)}
 \le\frac{2i}{(n+1)(n-1)}.
\]
Telescoping the products bounds their difference by the sum of the
factor differences. Taking the $m$th power multiplies the bound by
at most $m$, by the mean-value theorem on $[0,1]$. The difference at
index $j$ is therefore at most $ma^j j(j-1)/((n+1)(n-1))$.
Sum over $j$, use $\sum_{j\ge0}j(j-1)a^j=2a^2/(1-a)^3$, and divide
by $H_{n-1}\ge1+a$. Finally $H_0=1$ and $H_1=1+a$.
\end{proof}

\begin{proof}[Proof of Theorem~\ref{thm:sunbin}]
The hypothesis gives $a\le1/4$. Because $a^2/(1+a)$ and
$(1-a)^{-3}$ are increasing,
$2a^2/((1-a)^3(1+a))\le 32/135<1/3$. Since
$(n-1)/(2n-1)\ge1/3$ for $n\ge2$, the lemma gives
\[
 \frac{H_n}{H_{n-1}}<
 1+\frac{m}{(n+1)(2n-1)}
 \le\left(1+\frac{1}{(n+1)(2n-1)}\right)^m
 =\left(\frac{B_{n-1}B_{n+1}}{B_n^2}\right)^m.
\]
The second inequality is Bernoulli's inequality. At $n=1$ the same
strict comparison follows from $H_1/H_0=1+a\le5/4<1+m/2\le(3/2)^m$.
Multiplication by $H_{n+1}/H_n\ge1$ now gives the result because the
factors $q^n$ cancel.
\end{proof}

\begin{proposition}\label{prop:sun-small-q}
For each real $m\ge48$ the choice $q=3^{-m}$ gives
$S_2^{(m)}(q)S_4^{(m)}(q)<S_3^{(m)}(q)^2$.
\end{proposition}
\begin{proof}
Let $M_r=\max_{0\le k\le r}\binom{r+k}{k}/3^k$.
The ratio of consecutive terms is $(r+k+1)/(3(k+1))$, so that
$M_2=1$, $M_3=4/3$, and $M_4=5/3$. By bounding a sum by its number
of terms times its maximum and retaining one term at $r=3$ we obtain
\[
 \frac{S_2^{(m)}(q) S_4^{(m)}(q)}{S_3^{(m)}(q)^2}
 \le 15\left(\frac{M_2M_4}{M_3^2}\right)^m
 =15(15/16)^m<1.
\]
Indeed $(16/15)^{12}\ge1+12/15+66/225>2$.
\end{proof}
\section{Relation to the corpus and the priority record}\label{sec:status}
We used the repository snapshot \texttt{5887cde} of 15 November 2026
\cite{vmrc}. We screened the catalog of 190 TeX files and searched their text for source identifiers and subject terms. Two specific transfers can be checked in the repository.
Manuscript [179], in the proof of its harmonic defect transfer, separates
a baseline degree score from a nonnegative remainder and rewrites an edge
sum by counting its endpoints. Lemma~\ref{lem:rank} instead starts with
ranks of vertex pairs and produces a signed linear edge score whose
optimizer is explicit. Manuscript [043] bounds a product by telescoping
one changed factor at a time in its counting-bound lemma and sums
length-dependent losses by differentiating a geometric series in its
two-sided operator comparison. Lemma~\ref{lem:sunbin-drift} adapts these
steps to an endpoint-normalized binomial sum. No graph or arithmetic
claim from a draft in the corpus is needed as an unproved premise here.

\paragraph{Graph extremality.}
The original graph paper records a proof over a smaller exponent range
and a computational check through $n=11$ at $1/n$ \cite{ksfd}. The
linear objective, sorted degrees and antidiagonal graph were already
studied in \cite{abd}; fixed-size variants of the linear problem are
also known \cite{adr}. The two-point contact in
\eqref{eq:rankgap} gives simultaneous optimality for all nondecreasing concave
weights and a sharp loss for repeated degrees. The 2026 chemical
paper proves the case $n\equiv0\pmod4$ at $1/n$ and also gives
counterexamples to the separate constant-exponent chemical conjecture
\cite{dfks}. Its divisible-by-four theorem allows a larger
constant in the numerator of $p=c/n$ than our all-residue bound.
The parity argument in Section~\ref{sec:chem} supplies the other
three classes.

\paragraph{Chromatic and binomial coefficients.}
Huh's first-iterate theorem and the McNamara--Sagan criterion predate
this work. The additional input in Section~\ref{sec:cycles} is the
local third-iterate obstruction and its explicit sign along cycles.
The BIRS report also states the precise question
\cite{birs}. For the $R_n$ sequence, the generating function and
recurrence are inputs to an all-index comparison. The exact polynomial
enclosures implement earlier ideas about ratio inequalities for
P-recursive sequences \cite{bysun,houzhang}. For Zhi-Wei Sun's sum,
the coefficientwise claim has a published online counterexample
\cite{zwsun}; the author separately formulated the scalar question
at $q=1$ after that counterexample.

\paragraph{Scope of the historical claims.}
For G1 and G2, each target is specified in the listed source and is
answered at its original quantifiers. For G3 the source papers,
the authors' follow-up and the workshop report document the original
status and the partial results described above. Searches using problem
numbers, author names, DOIs, sequence identifiers, degree differences,
higher iterates, and equivalent inequalities did not identify prior
full solutions to the remaining targets. For G4 the proofs give either
all allowed parameters or an admissible counterexample. For G5 the
rank support bound, the parity estimate, the signed quartic, and
the quantitative binomial bounds provide substantive additional
conclusions. For G6, priority and the historical novelty of a general
method require independent comparison. The mathematical statements
above carry no assumption of historical priority or peer review.

\appendix
\section{Exact certificate for the ratio recurrence}\label{app:ratio}
The following Python code uses SymPy polynomial arithmetic in
$\mathbb Q(\sqrt2)[h]$. All sign tests compare rational numbers or their
squares. The degree checks cover every coefficient position, including
a zero if one were present. Together with the interval induction in
Section~\ref{sec:sun}, they certify the unbounded range.
\begin{lstlisting}[language=Python]
from sympy import QQ, sqrt
from sympy.polys import ring

K = QQ.algebraic_field(sqrt(2))
s = K.from_sympy(sqrt(2))
lam, t = 3 + 2*s, 3 - 2*s
ring_h, h = ring('h', K)
b = (84 - 39*s)/32
c = (-537 + 273*s)/64
d = (108996 - 83715*s)/4096

def H(x, sign):
    return x**5 - (K(3)/2)*x**4 + b*x**3 + c*x**2 + d*x + sign*20

def positive(q):
    # In K the internal coefficients are in descending powers of sqrt(2).
    coeff = q.to_list()
    if len(coeff) <= 1:
        return bool(coeff and coeff[0] > 0)
    u, v = coeff                 # u*sqrt(2) + v
    if u*v >= 0:
        return (u > 0 or v > 0)
    if v > 0:
        return v*v > 2*u*u
    return 2*u*u > v*v

x = 24 + h
D = [(x+i)**5 for i in range(4)]
l = [H(x+i, -1) for i in range(4)]
u = [H(x+i, 1) for i in range(4)]
def residual(a, b, c):
    return ((7*x+13)*t*a*b*D[2] - (7*x+15)*t**2*a*D[1]*D[2]
            + (x+1)*t**3*D[0]*D[1]*D[2] - (x+3)*a*b*c)
A = l[1]*D[0] - u[0]*D[1]
B = l[3]*D[2] - u[2]*D[3]
E = u[2]*D[1] - l[1]*D[2]
tests = [l[0] - D[0]/2, residual(u[0],l[1],l[2]),
         -residual(l[0],u[1],u[2]), l[1]**2*D[0]*D[2]-u[0]*u[2]*D[1]**2,
         A, B, l[2]*A*B*D[1] - u[0]*E**2*D[3]]
for p, degree in zip(tests, [5,11,11,17,8,8,24]):
    assert p.degree() == degree
    assert all(positive(p.get((j,), K.zero)) for j in range(degree+1))

R = [-1, 1, 7]
for n in range(26):
    q, rem = divmod((7*n+13)*R[-1]-(7*n+15)*R[-2]+(n+1)*R[-3], n+3)
    assert rem == 0
    R.append(q)
for n in (24,25):
    m = K(n)
    v = m**5 * R[n+1]
    assert positive(v-lam*R[n]*H(m,-1))
    assert positive(lam*R[n]*H(m,1)-v)
def ratio_delta(n):
    return R[n+1]**3*R[n-1]-R[n]**3*R[n+2]
def V(n):
    return R[n-1]*R[n+1]-R[n]**2
def minor_delta(n):
    return V(n-1)*V(n+1)-V(n)**2
assert ratio_delta(4) == -4623352
assert all(ratio_delta(n)>0 for n in range(5,25))
assert minor_delta(5) == -16234176
assert all(minor_delta(n)>0 for n in range(6,26))
assert minor_delta(7) == 343413129984
print("All exact checks passed.")
\end{lstlisting}
\begin{thebibliography}{99}
\bibitem{ksfd} M. Knor, R. \v{S}krekovski, S. Filipovski and D. Dimitrov,
Extremizing antiregular graphs by modifying total $\sigma$-irregularity.
\emph{Appl. Math. Comput.} \textbf{490} (2025), 129199.
Preprint 3 November 2024. \href{https://arxiv.org/html/2411.01530v1}{Preprint},
\href{https://doi.org/10.1016/j.amc.2024.129199}{journal}.
\bibitem{dfks} D. Dimitrov, S. Filipovski, M. Knor and R. \v{S}krekovski,
Modified total sigma-irregularity in chemical graphs.
\emph{Discrete Math. Lett.} \textbf{18} (2026), 1--7.
\href{https://doi.org/10.47443/dml.2026.041}{doi:10.47443/dml.2026.041}.
\bibitem{abd} H. Abdo, S. Brandt and D. Dimitrov, The total irregularity
of a graph. \emph{Discrete Math. Theor. Comput. Sci.} \textbf{16} (2014),
201--206. \href{https://doi.org/10.46298/dmtcs.1263}{Article}.
\bibitem{adr} H. Abdo, D. Dimitrov and G. Rote, Graphs of prescribed order
and size with maximal total irregularity. Manuscript, November 2025.
\href{https://page.mi.fu-berlin.de/rote/Papers/abstract/Graphs+of+prescribed+order+and+size+with+maximal+total+irregularity.html}{Author's record}.
\bibitem{am} T. Amdeberhan and V. H. Moll, Infinite log-convexity.
\emph{Online J. Anal. Comb.} \textbf{17}, paper 6, 1--10, Conjecture 21.
\href{https://combinatorialpress.com/article/ojac/vol17/305.pdf}{Journal PDF}.
\bibitem{birs} M. B\'ona, R. G\'omez, M. D. Ward and C. Mailler,
\emph{Analytic and Probabilistic Combinatorics}, BIRS workshop 22w5004,
13--18 November 2022, open problem session, Problem~5, p.~6.
\href{https://www.birs.ca/workshops/2022/22w5004/report22w5004.pdf}{Workshop report}.
\bibitem{huh} J. Huh, Milnor numbers of projective hypersurfaces and
the chromatic polynomial of graphs. \emph{J. Amer. Math. Soc.} \textbf{25}
(2012), 907--927. \href{https://doi.org/10.1090/S0894-0347-2012-00731-0}{Article}.
\bibitem{mcs} P. R. W. McNamara and B. E. Sagan, Infinite log-concavity:
developments and conjectures. \emph{Adv. Appl. Math.} \textbf{44} (2010),
1--15. \href{https://arxiv.org/abs/0808.1065}{Preprint}.
\bibitem{bysun} B. Y. Sun, On a ratio monotonicity conjecture of a new kind
of numbers. \emph{J. Inequal. Appl.} (2018), paper 24, Conjectures 4.7--4.8.
\href{https://arxiv.org/abs/1512.01008}{2015 preprint},
\href{https://doi.org/10.1186/s13660-018-1614-1}{journal}.
\bibitem{houzhang} Q.-H. Hou and Z.-R. Zhang, Asymptotic $r$-log-convexity
and P-recursive sequences. \href{https://arxiv.org/abs/1609.07840}{arXiv:1609.07840}.
\bibitem{zwsun} Z.-W. Sun, On the polynomials
$\sum_{k=0}^n\binom{n+k}{k}^m q^k$. MathOverflow question 479301,
21 September 2024; answer of Q. Zhou, 24 July 2026; scalar conjecture in
Sun's comment of 25 July 2026.
\href{https://mathoverflow.net/questions/479301/}{Original question, answer and comment}.
\bibitem{vmrc} QEEDDD, \emph{Verifiable Mathematical Reasoning Corpus},
manuscripts [043] and [179], snapshot \texttt{5887cde}.
\href{https://github.com/QEEDDD/Verifiable-Mathematical-Reasoning-Corpus/tree/5887cdea4898f0a3aa5bf7a8e0829697bd2b0f4b}{Versioned repository}.
\end{thebibliography}
\end{document}
